% Discussion for lab3_strace -- never overwritten by make_latex.py.

\texttt{strace} prints each system call a program makes. After loading its
libraries, \texttt{cat} opens \texttt{foo} with \texttt{openat}, which returns
file descriptor~\straceReadFd, and then copies it in four calls:
\begin{enumerate}
  \item \texttt{read(\straceReadFd, "hello\textbackslash n", \straceBufSize) = \straceBytes}:
    asks for up to \straceBufSize{}~bytes from \texttt{foo} and gets all
    \straceBytes.
  \item \texttt{write(\straceWriteFd, "hello\textbackslash n", \straceBytes) = \straceBytes}:
    writes those bytes to standard output.
  \item \texttt{read(\straceReadFd, "", \straceBufSize) = \straceEofRet}:
    returns~\straceEofRet, meaning end of file.
  \item \texttt{close(\straceReadFd) = \straceCloseRet}: releases the
    descriptor.
\end{enumerate}

So \texttt{3} is the descriptor of \texttt{foo}, \texttt{1} is standard output,
\texttt{6} is the number of bytes in \texttt{foo} (\texttt{hello} plus a
newline), and \texttt{32768} is the buffer size \texttt{cat} asks
\texttt{read()} to fill. GNU \texttt{cat}~9.7 asks for \straceBufSize{}~bytes
(\straceBufKiB\,KiB) instead; the role is the same. Ubuntu's default
\texttt{cat} (uutils) also opens \texttt{foo} on fd~\straceReadFd{} but copies
it inside the kernel with \texttt{splice()}, so it makes no
\texttt{read}/\texttt{write} calls at all.
